\section{Authorization-Aware Pre-Commit Prototype}

\begin{figure}[t]
\centering
\begin{tikzpicture}[
  node distance=5mm,
  box/.style={draw, rounded corners=1.5pt, align=center, font=\small, minimum height=8mm, fill=gray!8, text width=27mm},
  arrow/.style={-Latex, thick}
]
\node[box] (intercept) {Intercept candidate call};
\node[box, right=of intercept] (expand) {Expand into atoms};
\node[box, right=of expand] (alias) {Canonicalize aliases};
\node[box, right=of alias] (auth) {Check authorization context};
\node[box, below=of auth] (prov) {Apply provenance overlay};
\node[box, left=of prov, fill=green!8] (decision) {ALLOW / DENY / ABSTAIN};
\draw[arrow] (intercept) -- (expand);
\draw[arrow] (expand) -- (alias);
\draw[arrow] (alias) -- (auth);
\draw[arrow] (auth) -- (prov);
\draw[arrow] (prov) -- (decision);
\end{tikzpicture}
\caption{Authorization-aware local pre-commit prototype pipeline. The prototype is a local mock contract for pre-commit mediation, not a complete permission system.}
\label{fig:authz-pipeline}
\end{figure}


The E55 local prototype tests a stronger interface: rather than inferring authorization only from generic text, the environment provides a typed authorization context and the mediator expands candidate calls into atoms before decision. This is a local mock contract. It should be read as feasibility evidence for explicit pre-commit authorization infrastructure, not as a deployable permission system.

\paragraph{Pipeline.}
The prototype performs five steps. First, it intercepts the candidate tool call before commit. Second, it expands the call into effect-resource-operation atoms. Third, it canonicalizes resource identifiers and aliases. Fourth, it checks each atom against the authorization context. Fifth, it applies provenance/control-source checks and emits \textsc{Allow}, \textsc{Deny}, or \textsc{Abstain}.

\paragraph{Domain coverage.}
The local contract covers five domains: email, calendar, file, Slack, and transaction. Each domain contains 120 rows in E55-v2, for 600 rows total. The stress axes include multi-resource expansion, operation shifts, provenance shifts, resource shifts, and same-effect surface changes. The domain examples are intentionally compact but cover security-relevant distinctions that coarse tool monitors often collapse.

\paragraph{Abstention semantics.}
Abstention is a first-class decision. In E55-v2, unknown resources, unresolved aliases, insufficient evidence fallback, or ambiguous draft/commit status are not silently allowed. The high-level goal is to convert dangerous automatic decisions into either denied or explicitly abstained pre-commit states.

\paragraph{Why this matters.}
The E50 bottleneck shows that tuple-level binding without a local authorization contract can make many decisions while missing unsafe resource/authorization cases. The E55 prototype asks the constructive follow-up: if the environment exposes the missing authorization interface and the mediator reasons over atoms, does the bottleneck reduce in a controlled setting? The answer is yes for the E55-v2 local contract, subject to the limitations and audit findings reported below.
